\documentclass[11pt,a4paper]{article}
\usepackage[margin=25mm]{geometry}
\usepackage{amsmath,amssymb,amsthm,booktabs}
\setlength{\emergencystretch}{2em}
\newtheorem{theorem}{Theorem}\newtheorem{lemma}[theorem]{Lemma}
\newcommand{\SL}{\operatorname{SL}}\newcommand{\PSL}{\operatorname{PSL}}
\newcommand{\PGL}{\operatorname{PGL}}\newcommand{\GL}{\operatorname{GL}}
\newcommand{\Fp}{\mathbb F_p}\newcommand{\PG}{\operatorname{PG}}
\newcommand{\ip}{\mathrm{IP}^{+}}
\title{A parabolic subgroup and a conic in the rank-five construction}
\author{Human Compression project}\date{14 October 2026}
\begin{document}\maketitle

We give the subgroup and intersection arguments omitted from Section 4 of
Leemans--Vandenschrick \cite{source}. The four source matrices are unchanged.
The two signs require different geometry: the second end fixes a point for
one sign and preserves a conic for the other. We retain the classical
rotation correspondence, its recursive intersection criterion
\cite[Lemmas 2.2 and 2.3]{source}, and the description of the automorphisms
of $\PSL_3(F)$ as semilinear collineations and their compositions with
inverse transpose. No classification of subgroups of $\PSL_3(F)$ or
$\PGL_2(F)$ enters the proof.

Throughout, $p\geq5$ is prime, $w$ has order three, and $F=\Fp(w)$.
Thus $F=\Fp$ when $p\equiv1\pmod3$, and $F=\mathbb F_{p^2}$ when
$p\equiv-1\pmod3$. Put $q=|F|$ and $k\in\{1,-1\}$, and set
\[
 a=\operatorname{diag}(1,w,w^2),\quad
 b=\begin{pmatrix}-1&0&-k\\0&-w^2&-kw^2\\0&0&w\end{pmatrix},\quad
 c=\begin{pmatrix}1&k&0\\0&k&1\\0&-1&0\end{pmatrix},\quad
 d=\begin{pmatrix}0&1&0\\0&0&1\\1&0&0\end{pmatrix}.
\]
These matrices have determinant one. When discussing projective groups we
use the same letters for their images. Write $e_1=[1,0,0]^t$ and
$\mu_6=\langle-w^2\rangle\leq F^\times$.

\begin{theorem}
The marking $(a,b,c,d)$ generates $\PSL_3(F)$, satisfies the string
rotation relations and the full rotation intersection property, and has
type $[3,6,6,3]$ when $k=1$ and $[3,6,3,3]$ when $k=-1$.
If $F=\Fp$, it defines a chiral polytope. If $F=\mathbb F_{p^2}$ with
$p\equiv-1\pmod3$, it defines a directly regular polytope with rotation
group $\PSL_3(F)$; its full automorphism group is the extension by the
Frobenius automorphism, of order twice that of the rotation group.
\end{theorem}

\section{One projection and its translation kernel}
We first work in $\SL_3(F)$. Let $\widetilde A=\langle a,b,c\rangle$.
Every element has block form
\[
 \begin{pmatrix}\epsilon&v\\0&M\end{pmatrix},\qquad
 \epsilon\in\{1,-1\},\quad \det M=\epsilon.
\]
Projection to $M$ will be denoted by $\pi$. Put $t=k(2w+1)$; it is
nonzero because $(2w+1)^2=-3$. Direct multiplication gives
\[
 u=b^2a^{-1}=I+tE_{13}-tE_{23}.
\]
The lower block of $u$ is an upper unipotent matrix with parameter $-t$.
Conjugation by the lower block of $a$ multiplies that parameter by $w^2$.
Since $F=\Fp(w)$, additive combinations give every upper unipotent
matrix in $\SL_2(F)$. The lower block of $c$ is
\[
 \begin{pmatrix}k&1\\-1&0\end{pmatrix}
 =\begin{pmatrix}1&-k\\0&1\end{pmatrix}
   \begin{pmatrix}0&1\\-1&0\end{pmatrix}.
\]
Consequently the image contains the matrix interchanging the two root
groups, and then all of $\SL_2(F)$ by elementary elimination. The lower
block of $b$ has determinant $-1$, so
\[
                  \pi(\widetilde A)=\{M\in\GL_2(F):\det M=\pm1\}.
\]

The kernel is an additive subgroup $V$ of the row space $F^2$, normal
under the natural action of $\SL_2(F)$. It cannot be zero. If it were,
the unique lift of the central matrix $-I_2$ would be central in
$\widetilde A$ and have the form
\[
                   z=\begin{pmatrix}1&v\\0&-I_2\end{pmatrix}.
\]
Commutation with $a$ forces $v=0$, since $w,w^2\ne1$. But
$\operatorname{diag}(1,-1,-1)$ does not commute with $c$: its nonzero
entry $k$ in position $(1,2)$ changes sign. This is impossible in odd
characteristic. Thus $V\ne0$. The action of $\SL_2(F)$ is transitive
on nonzero row vectors, so normality forces $V=F^2$. We have proved
\begin{equation}\label{eq:A}
 \widetilde A=
 \left\{\begin{pmatrix}\epsilon&v\\0&M\end{pmatrix}:
 \epsilon=\pm1,\ \det M=\epsilon,\ v\in F^2\right\}.
\end{equation}
Its intersection with the scalar center of $\SL_3(F)$ is trivial.
The projective image $A$ therefore has order $2q^3(q^2-1)$.
Equivalently, $A$ consists of the projective transformations fixing
$e_1$ whose eigenvalue there, in an $\SL_3(F)$ representative, belongs
to $\mu_6$. Multiplication by a central cube root changes that
eigenvalue to $1$ or $-1$. In particular $A$ is generally smaller than
the full point stabilizer.

Equation \eqref{eq:A} supplies all $I+sE_{12}$ and $I+sE_{13}$.
Conjugation by $d$ cycles the coordinate indices and supplies the other
four elementary root groups. Hence
\[
                       \langle a,b,c,d\rangle=\SL_3(F)
\]
before taking projective images.

\section{The first adjacent pair}
Let $T=\{I+sE_{13}+tE_{23}:s,t\in F\}$ and let
$D_0=\operatorname{diag}(w,w^2)$ act on its two parameters. The vector
of $u=b^2a^{-1}$ is $t(1,-1)$. Its orbit spans exactly two dimensions
over $\Fp$: the vectors $t(1,-1)$ and $t(w,-w^2)$ are independent,
and $D_0^2+D_0+I=0$. Denote this additive subgroup by $V_0$.
Conjugation by $b$ acts on the parameters as $-D_0^2$, so preserves
$V_0$; conversely its square acts as $D_0$. Since $a=u^{-1}b^2$,
\begin{equation}\label{eq:H}
       H=\langle a,b\rangle=V_0\rtimes\langle b\rangle,
       \qquad |H|=6p^2.
\end{equation}
Here $b$ has order six, and its nonidentity powers have nontrivial
diagonal image. When $F=\Fp$, $V_0=T$; when $F=\mathbb F_{p^2}$,
it has order $q$, not $q^2$. All elements of $H$ fix both $e_1$ and
$e_2$.

For $1\leq j\leq5$ the entry $(b^j)_{13}$ is nonzero. Indeed it is
\[
                -k\,\frac{(-1)^j-w^j}{-1-w},
\]
and $-w^2$ has order six. Thus no such power is diagonal, proving
$\langle a\rangle\cap\langle b\rangle=1$ even projectively.

\section{The point case}
Suppose $k=1$ and put $r=[1,w,w^2]^t$.
Both $b,c$ fix $e_1$ and $r$. The source matrix
\[
 J(w)=\begin{pmatrix}w&w^2&1\\w^2&w&1\\1&1&1\end{pmatrix}
\]
has determinant $-3(2w+1)\ne0$. Literal multiplication gives
\begin{align*}
 J^{-1}a(w)J&=w^2d^{-1},&
 J^{-1}b(w)J&=w c(w^2)^{-1},\\
 J^{-1}c(w)J&=w^2 b(w^2)^{-1},&
 J^{-1}dJ&=w a(w^2)^{-1}.
\end{align*}
Thus it exchanges the marking with the inverse-reversed marking at
$w^2$, projectively. Applying \eqref{eq:A} at $w^2$ shows that
$B=\langle b,c,d\rangle$ is the corresponding restricted stabilizer
of $r$: its allowed eigenvalues at $r$ are precisely $\mu_6$.

Choose the basis $(e_1,r,e_2)$, represented by
\[
 P=\begin{pmatrix}1&1&0\\0&w&1\\0&w^2&0\end{pmatrix}.
\]
In that basis the two middle generators are
\[
 b'=\operatorname{diag}(-1,w,-w^2),\qquad
 c'=\begin{pmatrix}1&0&-w^2\\0&-w^2&-w\\0&0&-w\end{pmatrix}.
\]
Their commutator $b'c'b'^{-1}c'^{-1}$ has translation vector
$(-(2w+1),w)$ in positions $(1,3),(2,3)$. Conjugation by $b'$ acts
on these two coordinates with multipliers $w$ and $-w^2$.

The prime-field minimal polynomials of these multipliers are distinct
and relatively prime. In the prime-field case this follows from their
different values; in the quadratic case their roots have orders three
and six, respectively. Polynomial projections in this action therefore
separate the two coordinates. Each nonzero coordinate then generates
all of $F$ under its multiplier, since $F=\Fp(w)$. Consequently
$E=\langle b,c\rangle$ contains the full two-dimensional translation
group in this basis.

Modulo that group, the two eigenvalue pairs at $(e_1,r)$ are
$(-1,w)$ and $(1,-w^2)$. They generate all
$\{\pm1\}\times\mu_6$: the second pair has order six and cannot
supply the first coordinate $-1$. They therefore give all twelve
allowed diagonal classes. After choosing the unique scalar adjustment
making the first eigenvalue $\pm1$, these are exactly the elements
of $A\cap B$. We conclude that
\begin{equation}\label{eq:pointintersection}
                      A\cap B=E,\qquad |E|=12q^2.
\end{equation}

The action of $V_0$ on $r$, whose third coordinate is nonzero, is free.
Thus the stabilizer of $r$ in $H$ has order at most six. It contains
$\langle b\rangle$, so $H\cap E=\langle b\rangle$.
Also $\langle b\rangle\cap\langle c\rangle=1$: in the displayed
basis $b'$ is diagonal, whereas the $(1,3)$ entry of $(c')^j$ is nonzero
for $1\leq j\leq5$, because its two diagonal values have ratio of
order six. These statements, \eqref{eq:H}, and the rotation criterion
prove the full intersection property for $(a,b,c)$.

The same conclusion at $w^2$, transported by $J$, proves the
intersection property for $(b,c,d)$. In particular
$E\cap\langle c,d\rangle=\langle c\rangle$. Combining this with
\eqref{eq:pointintersection} gives
\[
             A\cap\langle c,d\rangle=\langle c\rangle.
\]
Finally $A\cap\langle d\rangle=1$, since $d$ cycles the three
coordinate points and $A$ fixes $e_1$.

\section{The conic case}
Suppose $k=-1$. The conic in the source is
\[
                       xy+w yz+w^2zx=0.
\]
Set $U=x+wz$, $V=z$, $W=y+w^2z$. Its equation becomes $UW-V^2=0$,
with parametrization $[s:t]\mapsto[s^2:st:t^2]$.
For $g=\begin{pmatrix}\alpha&\beta\\\gamma&\delta\end{pmatrix}$,
define
\[
 \iota(g)=\frac1{\det g}
 \begin{pmatrix}
 \alpha^2&2\alpha\beta&\beta^2\\
 \alpha\gamma&\alpha\delta+\beta\gamma&\beta\delta\\
 \gamma^2&2\gamma\delta&\delta^2
 \end{pmatrix}.
\]
This is the symmetric-square representation divided by its determinant;
it is a homomorphism into $\SL_3(F)$ and is unchanged by scalar
multiplication of $g$. Its projective action is faithful on
$\PGL_2(F)$: an element fixing every quadratic point fixes every
point of the projective line and is scalar.

Let $T$ denote the above coordinate change. Exact multiplication gives
\[
 w^2TbT^{-1}=\iota\begin{pmatrix}-w^2&0\\0&1\end{pmatrix},\quad
 wTcT^{-1}=\iota\begin{pmatrix}w&-w^2\\0&1\end{pmatrix},\quad
 w^2TdT^{-1}=\iota\begin{pmatrix}-w^2&1\\-w&0\end{pmatrix}.
\]
The scalar factors have order three and disappear projectively. We may
therefore work with these three $2\times2$ matrices, denoted by
$\bar b,\bar c,\bar d$.

The first two act on an affine coordinate as
\[
                        z\mapsto-w^2z,\qquad z\mapsto wz-w^2.
\]
Their commutator is a nonzero translation. Conjugation by their
multipliers gives its entire $F$-span over $\Fp$, hence all
translations, while the multipliers generate $\mu_6$. Thus
\begin{equation}\label{eq:affine}
 \langle\bar b,\bar c\rangle
 =\{z\mapsto tz+h:t\in\mu_6,\ h\in F\},\qquad |E|=6q.
\end{equation}
The element $\bar d$ sends infinity to $w$. Conjugating the full
translation group and then translating $w$ to zero gives the full
opposite unipotent root group. The two root groups generate
$\PSL_2(F)$. The determinants of $\bar c,\bar d$ are $w$, a square;
the determinant of $\bar b$ is $-w^2$. It follows directly from the
determinant-square quotient that
\[
 B\cong\begin{cases}
 \PSL_2(F),&-1\text{ is a square in }F,\\
 \PGL_2(F),&-1\text{ is not a square in }F.
 \end{cases}
\]
No subgroup classification is needed for this conclusion.

Under the coordinate change, $e_1$ is the conic point infinity and
$e_2$ is the point zero. An element of $B$ fixing $e_1$ is represented
by $\begin{pmatrix}t&h\\0&1\end{pmatrix}$. Its eigenvalue at $e_1$
in the displayed determinant-one lift is $t$. By \eqref{eq:A}, it
belongs to $A$ exactly when $t\in\mu_6$. These multipliers all occur
in \eqref{eq:affine}, so once again
\[
                              A\cap B=E.
\]
Since $H$ also fixes $e_2$, its intersection with $E$ fixes zero in
the affine model. This subgroup is exactly $\langle\bar b\rangle$;
therefore $H\cap E=\langle b\rangle$. The affine transformation
$\bar c$ has order three, and neither of its nonidentity powers fixes
zero. Hence $\langle b\rangle\cap\langle c\rangle=1$. Together
with the first cyclic intersection, this proves $\ip$ for $(a,b,c)$.

The matrices $c,d$ permute the four points
\[
                  e_1,e_2,e_3,[1,1,1]^t
\]
as two three-cycles: $c$ fixes $e_1$, and $d$ cycles $e_1,e_2,e_3$.
These points form a projective frame, so the action is faithful and
$\langle c,d\rangle\cong A_4$. Its stabilizer of $e_1$ is
$\langle c\rangle$, giving
\[
 A\cap\langle c,d\rangle=\langle c\rangle,
 \qquad A\cap\langle d\rangle=1.
\]
This supplies the actual conic intersection argument, rather than the
point-pair shorthand printed at the end of source Section 4.

\section{The full intersection property and orientation}
For either sign we have proved $\ip$ for $(a,b,c)$ and all three
recursive intersections
\[
 A\cap\langle b,c,d\rangle=\langle b,c\rangle,\quad
 A\cap\langle c,d\rangle=\langle c\rangle,\quad
 A\cap\langle d\rangle=1.
\]
The retained recursive criterion therefore gives the full rank-five
rotation intersection property. The six interval products are
nonidentity projective involutions. Literal multiplication gives squares
$I$ for $ab,bc,cd,abc,abcd$, and $(bcd)^2=w^2I$. Their nontriviality
is apparent from the nonzero off-diagonal entries. The eigenvalues of
$a,b,c,d$ give their projective orders $3,6,6,3$ for $k=1$ and
$3,6,3,3$ for $k=-1$.

Suppose first that $F=\Fp$. An automorphism reversing the marking must
invert $a$ and fix the three prefix involutions $ab,abc,abcd$. It
preserves $A=\langle a,ab,abc\rangle$. This group fixes $e_1$ and
fixes no line: its translation subgroup rules out lines not containing
$e_1$, and its lower $\SL_2(F)$ rules out a fixed line through $e_1$.
An automorphism involving inverse transpose would send its fixed point
to a fixed line, so cannot preserve $A$. There is no nontrivial field
automorphism. The proposed reverser is therefore conjugation by a
matrix $P\in\GL_3(F)$.

Projective commutation with a determinant-one involution whose square is
$I$ is ordinary commutation: its scalar multiplier has both square and
cube one, so equals one. Here
\[
 ab=\begin{pmatrix}-1&0&-k\\0&-1&-k\\0&0&1\end{pmatrix},\quad
 abc=-P_{23},\qquad abcd=-P_{12},
\]
where $P_{ij}$ is the coordinate transposition matrix. Commuting with
$P_{12},P_{23}$ forces $P=\alpha I+\beta\mathcal J$, with
$\mathcal J$ the all-ones matrix. The $(3,1)$ entry of
$\mathcal J(ab)-(ab)\mathcal J$ is $-2$, so $\beta=0$.
Thus $P$ is scalar and cannot invert the element $a$ of order three.
This proves chirality over the prime field for both signs.

If $F=\mathbb F_{p^2}$ with $p\equiv-1\pmod3$, Frobenius sends
$w$ to $w^2$ and hence $a$ to $a^{-1}$. All three displayed prefix
involutions are defined over $\Fp$, so Frobenius fixes them. It is an
involutory automorphism of the rotation group and is outer. The retained
correspondence gives the directly regular extension by this automorphism;
its full group, rather than its rotation subgroup, has order
$2|\PSL_3(F)|$.

\section{The six marking classes and attribution}
Over $\Fp$, the two choices $w,w^2$ give different ordered marking
classes for each sign. Indeed they are exactly the enantiomorphic
markings: $a(w^2)=a(w)^{-1}$, $b(w^2)=a(w)^2b(w)$, and $c,d$ are
unchanged. An isomorphism of these ordered markings would be the
reverser just excluded. For $k=1$ the matrix $J$ shows that their two
dual markings give the same two classes. For $k=-1$ the dual type is
$[3,3,6,3]$, different from both original types; it contributes two
further classes. Thus the source construction gives six ordered,
oriented marking classes, including duals.

The word ``polytopes'' in the source's count must be read with this
orientation convention. The two enantiomorphic markings are the two
choices of base-flag orbit in one unbased chiral polytope. If orientation
is forgotten, the above six classes give three unbased isomorphism
classes: one of type $[3,6,6,3]$, and a dual pair of the other two types.
The first is self-dual after orientation is forgotten. These distinctions
follow from the retained rotation correspondence; they do not alter
the source's existence theorem.

The matrices, the two signs, the conic, the duality matrix, and the
prime and quadratic-field constructions are the source authors' work.
The projection-kernel argument, the separated translation modules, the
explicit symmetric-square model and the full recursive intersections
above provide the omitted uniform proofs. We also make the full regular
group and the orientation convention explicit. This component treats
Theorem 1.2 and Section 4 only. It does not settle the source's later
classifications, global rank obstruction, or open questions.

\begin{thebibliography}{2}
\bibitem{source} D. Leemans and A. Vandenschrick,
\emph{On chiral polytopes having a group $\PSL(3,q)$ as automorphism
group}, fixed 25-page author edition, 2022, Theorem 1.2 and Section 4.
\bibitem{SW} E. Schulte and A. I. Weiss, \emph{Chiral polytopes},
DIMACS 4 (1991), 493--516; rotation correspondence and recursive
intersection criterion, retained as in source Section 2.
\end{thebibliography}
\end{document}
